%=== CHAPTER THREE ===
%% Kinematics, dynamics, and the parameter structure of the payload-adaptive
%% model.  Rewritten 2026-09-18 against the deployed implementation, in which
%% the estimator augments the composite mass AND the horizontal first mass
%% moment s = m_P r_xy.  Notation macros live in main.tex.

\chapter{Modeling and Parameter Structure}
\label{ch:model}
\begin{spacing}{1.5}
\setlength{\parskip}{0.22in}

A quadrotor that picks up a parcel does not merely become heavier. It also
becomes harder to rotate, and its centre of mass no longer lies on the thrust
axis. These are three distinct changes to the equations of motion, and a
controller that compensates only the first will fail in ways that mass
adaptation cannot explain.

This chapter derives the model that the estimator and the controller share, and
then makes explicit which of its parameters are identified online and which are
not. The structural claim of the chapter --- and the reason the method needs so
little sensing --- is that the three payload-induced quantities are not three
independent unknowns, but neither are they functions of the mass alone. Two
scalars suffice: the composite mass $\mT$ and the horizontal \emph{first mass
moment} $\svec=\mP\rxy$. From this pair the centre-of-mass offset, the
trim moment and the dominant roll/pitch inertia increment all follow by
rigid-body algebra, without the estimator ever being told the grasp offset, and
without any quantity becoming undefined when the payload leaves.

Section~\ref{sec:model:frames} fixes frames and conventions,
Sections~\ref{sec:model:kin}--\ref{sec:model:dyn} derive the kinematics and
dynamics, and Section~\ref{sec:model:payload} derives the classical payload
parameterisation in terms of mass and a known lever arm.
Section~\ref{sec:model:moment} then shows why that parameterisation is the wrong
coordinate system for a vehicle that must also \emph{release} its load, and
replaces it with the first-moment parameterisation that the deployed estimator
uses. Sections~\ref{sec:model:phys}--\ref{sec:model:ctrl} describe how the
parameters are measured and consumed, and
Sections~\ref{sec:model:whynot}--\ref{sec:model:approx} defend the design and
declare its approximations.

\section{Frames, Notation, and Conventions}\label{sec:model:frames}

Two reference frames are used throughout. The world frame $\mathcal{W}$ is
East--North--Up (ENU), so that gravity acts along $-z$. The body frame
$\mathcal{B}$ is Front--Left--Up (FLU) with its origin at the geometric centre of
the rotor disc --- \emph{not} at the centre of mass. This choice is deliberate
and consequential: the thrust vector passes through the body origin by
construction, so any offset between that origin and the true centre of mass
appears explicitly as a moment (Section~\ref{sec:model:dyn}) rather than being
absorbed silently into the frame definition.

The state vector and the control input are
\begin{equation}
x = \bigl[\bm{p};\, \bm{v};\, \bm{q};\, \bm{\omega}\bigr] \in \mathbb{R}^{13},
\qquad
u = \bigl[T;\, \bm{\tau}\bigr] \in \mathbb{R}^{4},
\label{eq:state}
\end{equation}
where $\bm{p}$ and $\bm{v}$ are the position and velocity in $\mathcal{W}$,
$\bm{q}$ is the unit quaternion encoding the attitude of $\mathcal{B}$ relative
to $\mathcal{W}$ (Section~\ref{sec:model:quat}),
$\bm{\omega}$ is the body rate in $\mathcal{B}$, $T$ is the collective thrust
along the body $z$ axis, and $\bm{\tau}$ is the body torque. We abbreviate
$\ev = [0,0,1]^{\!\top}$, and write $\mB$, $\mP$ and $\mT=\mB+\mP$ for the
bare-airframe, payload and composite masses. A subscript $xy$ denotes the
horizontal part of a body-frame vector, so that $\rxy$ and $\cvec$ are the
horizontal parts of $\rp$ and $\rc$ respectively.

\subsection{Attitude representation}\label{sec:model:quat}

Attitude is represented by a unit quaternion under the \emph{Hamilton} convention
with the scalar part first,
\begin{equation}
\bm{q} = [q_w,\, q_x,\, q_y,\, q_z]^{\!\top},
\qquad \|\bm{q}\| = 1,
\label{eq:quatdef}
\end{equation}
encoding the rotation from $\mathcal{B}$ to $\mathcal{W}$. The corresponding
rotation matrix is
\begin{equation}
R(\bm{q}) =
\begin{bmatrix}
q_w^2 + q_x^2 - q_y^2 - q_z^2 & 2(q_xq_y - q_wq_z) & 2(q_xq_z + q_wq_y) \\[2pt]
2(q_xq_y + q_wq_z) & q_w^2 - q_x^2 + q_y^2 - q_z^2 & 2(q_yq_z - q_wq_x) \\[2pt]
2(q_xq_z - q_wq_y) & 2(q_yq_z + q_wq_x) & q_w^2 - q_x^2 - q_y^2 + q_z^2
\end{bmatrix}.
\label{eq:Rq}
\end{equation}

Three conventions are fixed here and used throughout: the Hamilton product, the
scalar-first component order of \eqref{eq:quatdef}, and $\bm{\omega}$ resolved
in $\mathcal{B}$. The last of these determines the side on which the body rate
enters \eqref{eq:quat}.

One scalar derived from \eqref{eq:Rq} is used repeatedly: the cosine of the tilt
angle between the body $z$ axis and the vertical,
\begin{equation}
\cos\theta \;=\; \ev^{\!\top} R(\bm{q})\, \ev \;=\; 1 - 2\bigl(q_x^2+q_y^2\bigr),
\label{eq:costheta}
\end{equation}
which is the fraction of the collective thrust that acts vertically.

\section{Rigid-Body Kinematics}\label{sec:model:kin}

Translational kinematics is the definition of velocity,
\begin{equation}
\dot{\bm{p}} = \bm{v}.
\label{eq:kin}
\end{equation}

Attitude kinematics follows from the quaternion product of the attitude with the
pure quaternion formed from the body rate. Writing $\bm{q} = (q_w, \bm{q}_v)$
with $\bm{q}_v = [q_x, q_y, q_z]^{\!\top}$, and using the Hamilton product
$(a_w, \bm{a}_v) \otimes (b_w, \bm{b}_v) =
(a_w b_w - \bm{a}_v \!\cdot\! \bm{b}_v,\;
a_w \bm{b}_v + b_w \bm{a}_v + \bm{a}_v \!\times\! \bm{b}_v)$, we obtain
\begin{equation}
\dot{\bm{q}} = \tfrac{1}{2}\, \bm{q} \otimes [0,\, \bm{\omega}]
= \tfrac{1}{2}
\begin{bmatrix}
-\bm{q}_v \cdot \bm{\omega} \\[2pt]
q_w \bm{\omega} + \bm{q}_v \times \bm{\omega}
\end{bmatrix}
= \tfrac{1}{2}\, \Xi(\bm{q})\, \bm{\omega},
\label{eq:quat}
\end{equation}
where the second equality expands the product and the third collects the result
into the $4 \times 3$ matrix
\begin{equation}
\Xi(\bm{q}) =
\begin{bmatrix}
-q_x & -q_y & -q_z \\
 q_w & -q_z &  q_y \\
 q_z &  q_w & -q_x \\
-q_y &  q_x &  q_w
\end{bmatrix}.
\label{eq:Xi}
\end{equation}
The implementation uses \eqref{eq:Xi}, a plain matrix--vector product that is
linear in $\bm{\omega}$. Note that the body rate enters on the \emph{right} of
the product in \eqref{eq:quat}. Had $\bm{\omega}$ been resolved in
$\mathcal{W}$, the correct expression would be
$\dot{\bm{q}} = \tfrac{1}{2} [0, \bm{\omega}] \otimes \bm{q}$ instead, differing
by a conjugation by $\bm{q}$.

The matrix form also exposes a property that the product form hides. For any
unit quaternion,
\begin{equation}
\Xi(\bm{q})^{\!\top}\Xi(\bm{q}) = I_3,
\qquad
\bm{q}^{\!\top}\Xi(\bm{q}) = \bm{0}^{\!\top},
\label{eq:Xiprops}
\end{equation}
as may be checked directly from \eqref{eq:Xi} using $\lVert\bm{q}\rVert=1$. The
second identity gives $\bm{q}^{\!\top}\dot{\bm{q}} = 0$, and therefore
\begin{equation}
\frac{\mathrm{d}}{\mathrm{d}t}\lVert\bm{q}\rVert^2
= 2\,\bm{q}^{\!\top}\dot{\bm{q}} = 0 :
\label{eq:normconserve}
\end{equation}
the unit-norm constraint is preserved by the kinematics themselves, rather than
maintained by repeated renormalisation. This matters inside the estimator, which
integrates \eqref{eq:quat} across the whole window before any measurement is
compared: a systematic drift of $\lVert\bm{q}\rVert$ would appear as an attitude
residual and be attributed to the payload parameters.

Neither \eqref{eq:kin} nor \eqref{eq:quat} contains a mass, an inertia, or any
other physical parameter: kinematics is pure geometry. \emph{The payload
therefore enters the model only through the dynamics}, which confines the entire
parameter discussion of this chapter to the three equations of
Section~\ref{sec:model:dyn}.

\section{Rigid-Body Dynamics}\label{sec:model:dyn}

\subsection{Translational dynamics}\label{sec:model:trans}

Newton's second law in $\mathcal{W}$, with thrust rotated from the body frame and
a linear aerodynamic drag term, gives
\begin{equation}
\dot{\bm{v}} = \frac{1}{\mT}\Bigl( R(\bm{q})\, \ev\, T - k_d\, \bm{v} \Bigr) - g\, \ev,
\label{eq:trans}
\end{equation}
where $\mT$ is the \emph{composite} mass of airframe plus payload, $k_d$ is a
fixed drag coefficient, and $g$ is the gravitational acceleration.

Equation \eqref{eq:trans} is where the mass --- and \emph{only} the mass ---
enters, and it enters in exactly one role: as the scalar gain $1/\mT$ from thrust
to acceleration. This single fact underwrites the observability argument of the
whole method. If a thrust value can be obtained independently of the model, then
the measured acceleration determines $\mT$ uniquely by \eqref{eq:trans}, and it
does so in ordinary near-hover flight, with no dedicated excitation manoeuvre.
The qualification ``independently of the model'' is essential and is the subject
of Section~\ref{sec:model:phys}.

It is worth recording the special case of \eqref{eq:trans} along the vertical,
because Section~\ref{sec:model:seed} uses it directly. Projecting \eqref{eq:trans}
onto $\ev$ and neglecting the vertical component of the drag term gives the
vertical force balance
\begin{equation}
\mT\bigl(g + a_z\bigr) \;=\; T\cos\theta ,
\label{eq:vbal}
\end{equation}
with $\cos\theta$ from \eqref{eq:costheta} and $a_z=\dot v_z$.

\subsection{Rotational dynamics}\label{sec:model:rot}

Let $\Jb = \mathrm{diag}(J_{xx}, J_{yy}, J_{zz})$ denote the inertia of the bare
airframe about its own centre of mass. We assume throughout that the airframe is
balanced about the body origin of Section~\ref{sec:model:frames}, so that the two
points coincide when no payload is attached. This is an assumption about the
vehicle, not a consequence of the frame choice, and it is load-bearing:
Section~\ref{sec:model:moment} relies on it to attribute the entire first mass
moment to the payload.

Let $\Jt$ denote the composite inertia, derived in Sections~\ref{sec:model:paxis}
and~\ref{sec:model:moment}. Euler's equation for a rigid body, written about the
composite centre of mass and in the body frame, is
\begin{equation}
\dot{\bm{\omega}} = \Jt^{-1}\Bigl( \bm{\tau} + \tcom(\cvec) - \bm{\omega} \times \Jt \bm{\omega} \Bigr).
\label{eq:rot}
\end{equation}
The term $\bm{\omega} \times \Jt \bm{\omega}$ is the usual gyroscopic coupling.
The term $\tcom$ is derived next.

That \eqref{eq:rot} is written \emph{about the composite centre of mass} is not a
stylistic choice. Euler's equation takes this form only when the inertia tensor
and the applied moments refer to the same point, and that point is either the
centre of mass or a point fixed in the inertial frame; a free-flying vehicle has
no such fixed point, so the centre of mass is the only admissible choice. Both
$\Jt$ and $\tcom$ must therefore be computed about the composite CoM --- which
\emph{moves} relative to the body frame when a payload is attached, and that
displacement is precisely why $\tcom$ exists. Mixing the two reference points is
a common and silent modelling error.

\subsection{The thrust eccentricity moment}\label{sec:model:tcom}

Thrust acts at the rotor-disc centre, which is the body origin. When a payload is
attached the composite centre of mass moves to some
$\rc = [c_x, c_y, c_z]^{\!\top} \neq \bm{0}$, and the thrust therefore no longer
passes through it. The moment about the CoM is the cross product of the position
of the application point \emph{relative to the CoM}, namely $-\rc$, with the
force $\ev T$:
\begin{align}
\tcom(\cvec) &= (-\rc) \times \ev T
= -\begin{bmatrix} c_x \\ c_y \\ c_z \end{bmatrix}
\times \begin{bmatrix} 0 \\ 0 \\ T \end{bmatrix} \nonumber \\[4pt]
&= -\begin{bmatrix} c_y T - 0 \\ 0 - c_x T \\ 0 - 0 \end{bmatrix}
= \begin{bmatrix} -c_y T \\ +c_x T \\ 0 \end{bmatrix}.
\label{eq:tcom}
\end{align}

Two properties of \eqref{eq:tcom} are used repeatedly in what follows.

The vertical offset $c_z$ does not appear. Its contribution to the cross product
vanishes identically, since a vertical displacement of the centre of mass is
collinear with the body-$z$ thrust; gravity acts at the centre of mass and
likewise exerts no moment about it. The offset relevant to the attitude dynamics
is therefore two-dimensional by construction rather than by truncation.

The moment $\tcom$ is non-zero in level hover and does not average to zero along a
trajectory. For the payload of Section~\ref{sec:model:magnitudes} it amounts to
roughly two thirds of the available roll authority, so station-keeping with an
eccentric load requires a persistent trim torque. A prediction model that omits
\eqref{eq:tcom} is therefore in error by a constant moment rather than by a small
one, and plans as though no trim torque were required.
Section~\ref{sec:model:ctrl} shows that the same term must additionally appear in
the input reference of the cost function.

\subsection{Structure of the parameter entry points}\label{sec:model:entry}

Collecting \eqref{eq:trans}, \eqref{eq:rot} and \eqref{eq:tcom}, the model has
the property on which the rest of this dissertation is built:

\begin{quote}
\emph{Each payload-dependent parameter has exactly one point of entry.} The
composite mass appears as $1/\mT$ in the translational dynamics; the inertia
increment appears only inside $\Jt$ in the rotational dynamics; the
centre-of-mass offset appears only in the eccentricity moment, which is also a
rotational quantity.
\end{quote}

This separation is what makes the parameters individually attributable to
distinct physical channels --- the mass to the thrust--acceleration relation, the
inertia to angular acceleration, the offset to the steady-state torque balance.
It is exploited directly in Sections~\ref{sec:model:mhe}
and~\ref{sec:model:whynot}.

\section{Payload-Induced Parameters}\label{sec:model:payload}

We now derive $\cvec$ and the inertia increment. The payload is modelled as a
point mass $\mP$ rigidly attached at a body-frame offset
$\rp = [r_x, r_y, r_z]^{\!\top}$ with $r_z < 0$ (the load hangs below the
airframe). Define the composite mass and the \emph{reduced mass}
\begin{equation}
\mT = \mB + \mP,
\qquad
\mu = \frac{\mB\, \mP}{\mB + \mP} = \frac{\mB\, \mP}{\mT}.
\label{eq:mu}
\end{equation}

This section treats $(\mP,\rp)$ as the parameterisation, because it is the
physically transparent one. Section~\ref{sec:model:moment} then shows that this
pair is not what an estimator can or should carry, and re-expresses every result
below in coordinates that survive release.

\subsection{Centre-of-mass offset}\label{sec:model:com}

The composite centre of mass is the mass-weighted mean of the two constituent
positions. Taking the airframe centre of mass at the body origin and the payload
at $\rp$,
\begin{equation}
\rc = \frac{\mB \cdot \bm{0} + \mP \cdot \rp}{\mT}
= \frac{\mP}{\mT}\, \rp ,
\label{eq:rc}
\end{equation}
whose horizontal components are the pair required by \eqref{eq:tcom}:
\begin{equation}
\cvec = [c_x,\, c_y]^{\!\top} = \frac{\mP}{\mT}\, [r_x,\, r_y]^{\!\top}.
\label{eq:cxy}
\end{equation}

\subsection{Inertia increment and the reduced mass}\label{sec:model:paxis}

The inertia increment requires more care than the parallel-axis theorem applied
naively, because \emph{the centre of mass itself moves when the payload is
attached}. The theorem must be applied about the \emph{new} centre of mass
\eqref{eq:rc}, and both bodies contribute, since the airframe no longer sits at
the CoM either.

Consider a single axis and let $d$ denote the corresponding lever arm. From
\eqref{eq:rc}, the distances of the two bodies from the composite centre of mass
are
\begin{equation}
\underbrace{\bigl\| \bm{0} - \rc \bigr\| = \frac{\mP}{\mT}\, d}_{\text{airframe}},
\qquad
\underbrace{\bigl\| \rp - \rc \bigr\| = \Bigl( 1 - \frac{\mP}{\mT} \Bigr) d = \frac{\mB}{\mT}\, d}_{\text{payload}} .
\label{eq:dists}
\end{equation}
Note that the payload's distance is \emph{smaller} than $d$. Adding the two
parallel-axis contributions and simplifying,
\begin{align}
\Delta J^{(d)}
&= \mB \Bigl( \frac{\mP}{\mT} d \Bigr)^{\!2}
 + \mP \Bigl( \frac{\mB}{\mT} d \Bigr)^{\!2} \nonumber \\[4pt]
&= \frac{d^2}{\mT^2}\Bigl( \mB\, \mP^2 + \mP\, \mB^2 \Bigr)
 = \frac{d^2}{\mT^2}\, \mB\, \mP \bigl( \mP + \mB \bigr) \nonumber \\[4pt]
&= \frac{\mB\, \mP}{\mT}\, d^2
 = \mu\, d^2 .
\label{eq:reduced}
\end{align}
The factor $(\mB + \mP)$ produced by the sum cancels one power of $\mT$ in the
denominator, and the two contributions collapse onto the reduced mass
\eqref{eq:mu} --- the same quantity that governs the two-body problem in
classical mechanics, and for the same reason: the inertia of a rigid pair depends
only on their \emph{relative} configuration.

Equation \eqref{eq:reduced} also identifies a natural error. Writing
$\Delta J^{(d)} = \mP d^2$ --- treating the payload as orbiting a fixed body
origin, and forgetting both the airframe's own contribution and the displacement
of the CoM --- overestimates the increment by the factor $\mT / \mB$, which is
\SI{14.5}{\percent} at the payload considered below. Two limiting cases confirm
\eqref{eq:reduced}: as $\mP \!\ll\! \mB$ we recover $\mu \to \mP$ (a light load
barely moves the CoM, and the naive expression becomes correct), while as
$\mP \!\gg\! \mB$ we obtain $\mu \to \mB$ (a heavy load anchors the CoM and it is
the airframe that revolves about it).

The same computation carried out for the full tensor rather than one axis gives
the increment in closed form. Writing $I_3$ for the identity,
\begin{align}
\Delta J(\mP,\rp)
&= \mB\Bigl(\lVert\rc\rVert^2 I_3 - \rc\rc^{\!\top}\Bigr)
 + \mP\Bigl(\lVert\rp-\rc\rVert^2 I_3 - (\rp-\rc)(\rp-\rc)^{\!\top}\Bigr) \nonumber\\[4pt]
&= \mu\Bigl(\lVert\rp\rVert^2 I_3 - \rp\,\rp^{\!\top}\Bigr),
\label{eq:dJfull}
\end{align}
by exactly the cancellation of \eqref{eq:reduced} applied entrywise. Written out,
\begin{equation}
\dJ = \mu
\begin{bmatrix}
r_y^2 + r_z^2 & -r_x r_y      & -r_x r_z \\[2pt]
-r_x r_y      & r_x^2 + r_z^2 & -r_y r_z \\[2pt]
-r_x r_z      & -r_y r_z      & r_x^2 + r_y^2
\end{bmatrix}.
\label{eq:dJmatrix}
\end{equation}
Each diagonal entry is \eqref{eq:reduced} evaluated with the lever arm of that
axis, and each off-diagonal entry is the product $-\mu r_i r_j$ of the two arms
it couples. Reading the diagonal off gives the familiar per-axis increments
\begin{equation}
\dJ_{xx} = \mu\bigl(r_y^2 + r_z^2\bigr), \quad
\dJ_{yy} = \mu\bigl(r_x^2 + r_z^2\bigr), \quad
\dJ_{zz} = \mu\bigl(r_x^2 + r_y^2\bigr).
\label{eq:djaxes}
\end{equation}
The composite inertia is $\Jt=\Jb+\dJ$.
Section~\ref{sec:model:Jofs} regroups the same six entries into the coordinates
that survive release.

\subsection{Magnitudes on the experimental platform}\label{sec:model:magnitudes}

To fix ideas, consider a \SI{0.3}{\kg} payload attached at
$\rp = [0,\, 0.109,\, -0.47]\,$\si{\meter} on the platform of this work
($\mB = \SI{2.0643}{\kg}$, $J_{xx} = \SI{0.0142}{\kg\meter\squared}$, roll/pitch
torque authority \SI{0.5}{\newton\meter}). Equations \eqref{eq:mu},
\eqref{eq:cxy} and \eqref{eq:djaxes} give
\begin{align}
\mu &= \SI{0.262}{\kg}, &
\dJ_{xx} &= \SI{0.0610}{\kg\meter\squared}, \nonumber \\
\frac{J_{xx} + \dJ_{xx}}{J_{xx}} &= 5.30, &
c_y &= \SI{13.8}{\milli\meter},
\label{eq:numbers}
\end{align}
and hence a constant hover trim torque
\begin{equation}
|c_y|\, \mT\, g = \SI{0.32}{\newton\meter} \approx \SI{64}{\percent}\ \text{of the roll authority.}
\label{eq:trimtorque}
\end{equation}

These numbers state the problem this dissertation addresses more sharply than any
qualitative argument. A \SI{0.3}{\kg} parcel on a \SI{2.06}{\kg} airframe is a
\SI{15}{\percent} change of mass --- but simultaneously a more than
\emph{five-fold} change of roll inertia and a standing demand on two thirds of
the torque margin. Adapting the mass alone accounts for neither of the latter
two. Note also that the $r_z^2$ term contributes \SI{97.4}{\percent} of
$\dJ_{xx}$: the increment is dominated by how far \emph{below} the airframe the
load hangs, not by how far it is off-centre. That asymmetry is exploited
repeatedly below, because $r_z$ turns out to be the one geometric quantity the
estimator cannot observe and, fortunately, the one whose error matters least.

\section{The First Mass Moment as the Payload State}\label{sec:model:moment}

Everything in Section~\ref{sec:model:payload} is expressed in the coordinates
$(\mP, \rp)$. Those are the natural coordinates for \emph{describing} a grasp.
They are the wrong coordinates for \emph{estimating} one, and the reason is
visible in \eqref{eq:cxy}: the map from $(\mP,\rp)$ to the observable quantities
is not invertible in the limit that matters most.

\subsection{Why the offset is the wrong coordinate}\label{sec:model:whyS}

Consider what each quantity does as the payload is released, $\mP \to 0$.

\begin{itemize}
\item The moment $\tcom$ of \eqref{eq:tcom} tends to zero, because
$\cvec = (\mP/\mT)\rxy \to \bm{0}$. This is physical and benign.
\item The horizontal offset $\rxy$ itself becomes \emph{undefined}. There
is no fact of the matter about where a payload of zero mass is attached, and no
measurement can supply one: every observable depends on $\rxy$ only
through the product $\mP\rxy$.
\end{itemize}

An estimator parameterised by $(\mP, \rxy)$ therefore carries a direction
that becomes unidentifiable exactly at release. If $\rxy$ is instead held
fixed as a prior, as in a conventional coupled model, the situation is worse
rather than better: the lever arm stays in the model after the load has gone, so
the model predicts a phantom trim moment $(\mP/\mT)\rxy{}T$ that the
measurements contradict. The only remaining way for the optimiser to cancel that
phantom is to drive $\mP$ --- and therefore $\mT$ --- toward or below the
bare-airframe value. This is not hypothetical: in the mass-only formulation that
preceded the present design, 14 of 55 post-release samples pinned at the mass
lower bound, and the recovered composite mass was biased low by
\SI{2.81}{\percent}. The information ``is there a payload at all'' had been
encoded into the mass channel, where it corrupted the quantity the controller
most needs.

The remedy is a change of coordinates. Define the horizontal \emph{first mass
moment} of the composite body about the body-frame origin,
\begin{equation}
\svec \;=\; \begin{bmatrix} s_x \\ s_y\end{bmatrix}
\;=\; \mP\,\rxy \quad [\si{\kg\meter}].
\label{eq:sdef}
\end{equation}
Because the bare airframe's centre of mass coincides with the body origin, the
airframe contributes nothing to the first moment and the entire moment is the
payload's, whatever grasp offset the gripper happens to realise. Equation
\eqref{eq:sdef} is the standard linear parameterisation of rigid-body inertial
parameters \cite{atkeson1986}: the equations of motion contain the \emph{product}
$\mP\rxy$ and never the two factors separately, which is precisely why the
product, and not the offset, is the quantity to estimate.

\subsection{The two channels the first moment drives}\label{sec:model:schannels}

Substituting \eqref{eq:sdef} into \eqref{eq:cxy} gives the first payoff:
\begin{equation}
\cvec \;=\; \frac{\mP}{\mT}\,\rxy
      \;=\; \frac{1}{\mT}\bigl(\mP\,\rxy\bigr)
      \;=\; \frac{\svec}{\mT}.
\label{eq:c_from_s}
\end{equation}
The centre-of-mass offset is available from the two estimated states alone. No
grasp geometry is required, and --- critically --- $\mP$ does not appear in a
denominator, so \eqref{eq:c_from_s} stays well defined and continuous as the
payload leaves. Combining \eqref{eq:c_from_s} with \eqref{eq:tcom}, the trim
moment consumed by the model is
\begin{equation}
\tcom \;=\; \frac{T}{\mT}\begin{bmatrix} -s_y \\ +s_x \\ 0\end{bmatrix}.
\label{eq:tcom_s}
\end{equation}

The second channel is the off-diagonal inertia. Taking the $(1,3)$ entry of
\eqref{eq:dJfull} and using \eqref{eq:mu} and \eqref{eq:sdef},
\begin{equation}
J_{xz} \;=\; -\mu\, r_x r_z
 \;=\; -\frac{\mB\,\mP}{\mT}\, r_x r_z
 \;=\; -\frac{\mB}{\mT}\,\bigl(\mP r_x\bigr)\, r_z
 \;=\; -\frac{\mB}{\mT}\, s_x\, r_z ,
\label{eq:Jxz}
\end{equation}
and identically $J_{yz} = -(\mB/\mT)\, s_y\, r_z$. The payload mass has cancelled
exactly: these terms need only $\svec$ and the vertical arm $r_z$.

\subsection{Re-expressing the inertia increment}\label{sec:model:Jofs}

It is convenient to split \eqref{eq:dJfull} into the three groups that behave
differently under the change of coordinates. Define
\begin{equation}
A = \mu\, r_z^2, \qquad
B_i = \mu\, r_i r_z = \frac{\mB}{\mT}\, s_i\, r_z, \qquad
C_{ij} = \mu\, r_i r_j \quad (i,j\in\{x,y\}),
\label{eq:ABC}
\end{equation}
so that the composite inertia of \eqref{eq:dJfull} reads
\begin{equation}
\Jt \;=\; \Jb \;+\; \bigl(A + C_{xx} + C_{yy}\bigr) I_3
\;-\;
\begin{bmatrix}
C_{xx} & C_{xy} & B_x \\
C_{xy} & C_{yy} & B_y \\
B_x    & B_y    & A
\end{bmatrix}.
\label{eq:JABC}
\end{equation}
Expanding the $(1,1)$ entry recovers $J_{xx}+\mu(r_y^2+r_z^2)$ and the $(3,3)$
entry recovers $J_{zz}+\mu(r_x^2+r_y^2)$, confirming \eqref{eq:djaxes}.

In the new coordinates the three groups have sharply different status.
\begin{enumerate}
\item $B_i$ requires \emph{only} $\svec$ and $r_z$, by \eqref{eq:Jxz}.
\item $A = (\mB\mP/\mT)\,r_z^2$ requires $\mP = \mT-\mB$ and $r_z$, both
available: the composite mass is estimated and $r_z$ is a weak prior. No
horizontal geometry is needed.
\item $C_{ij}$ is the problem. Writing it in the new coordinates,
\begin{equation}
C_{ij} \;=\; \mu\, r_i r_j
\;=\; \frac{\mB \mP}{\mT}\cdot\frac{s_i}{\mP}\cdot\frac{s_j}{\mP}
\;=\; \frac{\mB}{\mT\,\mP}\; s_i s_j .
\label{eq:Cdegenerate}
\end{equation}
\end{enumerate}

Equation \eqref{eq:Cdegenerate} is the formal statement of the difficulty. As
$\mP\to0$ with $\svec$ not yet driven to zero by the data --- exactly the state
of affairs during the first fraction of a second after release --- the
coefficient $\mB/(\mT\mP)$ diverges. The model then admits a region in which a
vanishingly light payload hangs on an arbitrarily long lever arm: physically
absurd, numerically attractive to an optimiser looking for a way to explain a
residual, and ill-conditioned in exactly the direction the estimator is trying to
resolve.

The deployed model therefore \emph{discards} the $C_{ij}$ group, setting
$C_{xx}=C_{yy}=C_{xy}=0$ in \eqref{eq:JABC}. What remains,
\begin{equation}
\Jt \;=\; \Jb + A\,I_3 -
\begin{bmatrix} 0 & 0 & B_x \\ 0 & 0 & B_y \\ B_x & B_y & A\end{bmatrix},
\qquad
\dJ_{xx}=\dJ_{yy}=A=\frac{\mB\,\mP}{\mT}\,r_z^2 ,
\label{eq:Jdeployed}
\end{equation}
contains no term that requires $\mP$ and $\rxy$ to be known separately.
The price is quantified exactly: the roll increment loses $C_{yy}=\mu r_y^2$ out
of $\mu(r_y^2+r_z^2)$, a relative error
\begin{equation}
\frac{C_{yy}}{A + C_{yy}} = \frac{r_y^2}{r_z^2+r_y^2},
\label{eq:Cerror}
\end{equation}
which at the platform geometry of \eqref{eq:numbers} ($r_y=\SI{0.109}{\meter}$,
$r_z=\SI{-0.47}{\meter}$) is \SI{2.6}{\percent}. That is the cost of removing a
degenerate direction from the model, and it buys a formulation in which the
post-release behaviour is well posed. Offline replay of a synthetic
grasp--release sequence makes the trade concrete: retaining $C_{ij}$ produced a
composite-mass estimate that climbed monotonically to \SI{3.17}{\kg} after
release and did not return, with $s_y$ changing sign frame to frame, whereas the
reduced model of \eqref{eq:Jdeployed} settled at \SI{-0.34}{\percent} of the bare
airframe with $\svec$ and $\dJ$ decaying together and no bound contact.

\subsection{The inertia lever and the frozen coefficient}\label{sec:model:Amode}

Removing $C_{ij}$ eliminates one pathological coupling between mass and geometry,
but not the only one. The remaining $A$ term is still linear in $\mP$, and that
is enough to give the optimiser a second lever on the mass channel.

The mechanism is as follows. Suppose a phantom horizontal offset $c_y$ survives
briefly after release. It generates a predicted roll acceleration of magnitude
\begin{equation}
\bigl|\dot\omega_x^{\text{ghost}}\bigr| \;\approx\; \frac{|c_y|\,T}{J_{xx} + A(\mP)},
\label{eq:ghost}
\end{equation}
which the measurements contradict, since the real vehicle is not rolling. The
numerator is bounded, because $\lVert\cvec\rVert\le\lVert\rxy\rVert$ is a
geometric limit. The denominator, however, grows linearly in $\mP$, and grows
fast: $r_z^2=\SI{0.22}{\meter\squared}$ is some fifteen times
$J_{xx}=\SI{0.0142}{\kg\meter\squared}$. \emph{Increasing the mass estimate
therefore dilutes the phantom moment.} The optimiser will do so if the rotational
residual it thereby removes outweighs the translational residual it thereby
creates --- and in offline replay of a \SI{0.3}{\kg} release it did exactly that,
driving $\hat m$ to \SI{3.17}{\kg} ($+\SI{53}{\percent}$) and accepting a
\SI{-3.4}{\meter\per\second\squared} translational inconsistency to suppress the
angular residual. The same disease has a second outlet: with the lower bound
active it instead pushes the estimate down onto $m_{\min}$. Either way, the mass
has been recruited as a geometry knob.

The implementation therefore evaluates $A$ from the previous window's converged
mass estimate $\hat m^{-}$, supplied to the current solve as a parameter rather
than computed from the decision variable:
\begin{equation}
A \;=\; \frac{\mB\,\bigl(\hat m^{-}-\mB\bigr)^{+}}{\hat m^{-}}\; r_z^2,
\qquad
\frac{\partial A}{\partial \mT}\Big|_{\text{window}} = 0 .
\label{eq:Afrozen}
\end{equation}
Within a single solve $A$ is then a constant and the gradient path of the
preceding paragraph is absent; between solves it is refreshed at the
\SI{10}{\hertz} window rate, so its value continues to follow the true mass with
one window of lag.

Table~\ref{tab:amode} lists the four treatments of $A$ that the implementation
supports and that Chapter~\ref{ch:evaluation} compares. They differ in two
independent respects: whether the derivative is broken, and whether the retained
value is correct. Only the lagged form satisfies both, and it is the deployed
default.

\begin{table}[ht]
\centering
\caption{Treatments of the inertia coefficient $A$. Breaking the derivative
removes the gradient path of \eqref{eq:ghost}; retaining a correct value is a
separate requirement, which the last two forms fail.}
\label{tab:amode}
\small
\begin{tabular}{@{}llll@{}}
\toprule
Treatment & $A$ evaluated from & $\partial A/\partial \mT$ & Retained value \\
\midrule
live & the decision variable $\mT$ & non-zero & exact \\
lagged (deployed) & previous solution $\hat m^{-}$ & zero & exact to one window \\
envelope & payload envelope $\mP^{\mathrm{env}}$ & zero & biased by $\mP^{\mathrm{env}}/\mP$ \\
omitted & --- ($A\equiv0$) & zero & zero; diagnostic only \\
\bottomrule
\end{tabular}
\\[4pt]
\footnotesize The corresponding implementation settings are listed in Appendix~A.
\end{table}

\subsection{Smooth positive part}\label{sec:model:smoothpos}

Two of the expressions above use $\mP=(\mT-\mB)^{+}$, and the mass decision
variable is permitted to fall slightly below $\mB$ to leave room for calibration
error. A hard $\max(\cdot,0)$ is not differentiable at the origin, which makes a
Gauss--Newton step chatter in precisely the regime --- $\mP$ near zero --- where
release decisions are taken. The implementation therefore uses the smooth
positive part
\begin{equation}
\mP^{+} \;=\; \tfrac{1}{2}\Bigl(\Delta + \sqrt{\Delta^2+\varepsilon^2}\Bigr),
\qquad \Delta = \mT-\mB, \quad \varepsilon=\SI{0.02}{\kg},
\label{eq:smoothpos}
\end{equation}
which satisfies $\mP^{+}\ge \tfrac{\varepsilon}{2}>0$, is smooth, and agrees with
$\max(\Delta,0)$ to within $\varepsilon/2$ everywhere. Positivity matters beyond
smoothness: a negative $\mP$ would make $\mu<0$ in \eqref{eq:ABC} and could
render $\Jt$ indefinite, at which point $\Jt^{-1}$ in \eqref{eq:rot} diverges.

\section{Model-Independent Reconstruction of the Physical Input}\label{sec:model:phys}

The thrust and torque appearing in \eqref{eq:trans} and \eqref{eq:rot} must be
known to the estimator. They are reconstructed from measured rotor speeds
$\omega_i$, the known rotor positions $(x_i,y_i)$ and the spin directions
$d_i\in\{\pm1\}$, via the quadratic rotor model $F_i = k_f \omega_i^2$:
\begin{equation}
\Tphys = \sum_i F_i, \qquad
\tauphys =
\begin{bmatrix}
\sum_i y_i F_i \\[2pt]
-\sum_i x_i F_i \\[2pt]
-k_m \sum_i d_i F_i
\end{bmatrix},
\label{eq:phys}
\end{equation}
with $|x_i|=|y_i|=\SI{0.174}{\meter}$,
$k_f=\SI{8.55e-6}{\newton\second\squared}$ and yaw-reaction coefficient
$k_m=\SI{0.016}{\meter}$. The first two moment components are lever-arm products;
the third is the aerodynamic reaction torque, proportional to rotor thrust and
opposed to the rotor's own spin.

\paragraph{This is a wrench, not a state estimate.} Equation \eqref{eq:phys}
returns the wrench the \emph{actuators generate}. It deliberately excludes
$\Jt\dot{\bm\omega}$ and rotor gyroscopic terms, which is exactly what makes it
usable as the known input $u^{\mathrm{phys}}$ of the estimator rather than as a
second, competing estimate of the state. Rotor speeds are sampled at
\SI{250}{\hertz} and averaged over each \SI{10}{\hertz} estimation interval,
which is the zero-order-hold-consistent discretisation of \eqref{eq:trans} and
\eqref{eq:rot}.

\paragraph{Why the commanded input cannot be used.} Let the controller compute
its thrust command from an assumed mass $\tilde{m}$, so that in hover
$T_{\mathrm{cmd}} = \tilde{m} g$. Feeding $T_{\mathrm{cmd}}$ into
\eqref{eq:trans} and solving for the mass yields
$\mT = T_{\mathrm{cmd}} / g = \tilde{m}$ --- for \emph{any} $\tilde{m}$
whatsoever. The relation is satisfied identically, the residual carries no
information, and the mass is unobservable: the estimate would simply reproduce
whatever the controller already believed. Any inversion error, saturation or
motor lag between command and rotor is likewise invisible. Equation
\eqref{eq:phys} breaks this circularity because rotor speeds are a physical
measurement into which no model state enters, taken downstream of allocation and
the motor dynamics. This is a recurring principle in what follows: an adaptive
loop must be fed by a signal that is not itself a product of the adaptation.

The yaw channel makes the point quantitatively. Across twenty batches of recorded
flights, the roll and pitch torques actually executed track the controller's
intended values closely (median correlation $+0.93$ and $+0.70$, slopes $1.08$
and $1.10$), but the executed yaw torque is only $0.32$ of the intended value: a
quadrotor's yaw authority comes solely from reaction torque and is an order of
magnitude weaker than its roll/pitch authority. An estimator fed the intended yaw
torque is fed a moment three times larger than the one the vehicle actually
experienced. Reconstructing all four components from \eqref{eq:phys} removes the
last place where the estimator could learn what the controller was thinking.

\paragraph{Thrust-to-throttle inversion and the actuator envelope.} The same
rotor model defines what the vehicle can actually produce, and therefore the
input constraint set of the controller. The flight stack accepts a normalised
throttle $\eta$, which the electronic speed controllers map affinely to rotor
speed before the quadratic thrust law applies:
\begin{equation}
\omega(\eta) = \omega_{\min} + \Delta\omega\,\eta,
\qquad
T(\eta) = K\,\omega(\eta)^2,
\qquad K = 4k_f ,
\label{eq:throttlemap}
\end{equation}
with $\omega_{\min}=\SI{150}{\radian\per\second}$ and
$\Delta\omega=\SI{850}{\radian\per\second}$. Inverting \eqref{eq:throttlemap}
gives the command actually published,
\begin{equation}
\eta(T) = \frac{\sqrt{T/K}-\omega_{\min}}{\Delta\omega},
\label{eq:throttleinv}
\end{equation}
which is exact at every mass, unlike the linear hover-referenced heuristic it
replaced. Because $\eta$ is clipped to $[0.05,\,0.95]$ --- the upper limit
reserving \SI{5}{\percent} of the command range so that the rate loop retains
differential authority when all four rotors are near saturation --- the
physically attainable collective thrust is
\begin{equation}
T \in \bigl[\,T(0.05),\; T(0.95)\,\bigr] = [\,\SI{1.27}{\newton},\; \SI{31.35}{\newton}\,].
\label{eq:envelope}
\end{equation}
Section~\ref{sec:model:ctrl} uses \eqref{eq:envelope} as the thrust box of the
optimal control problem. This matters more than it may appear: the earlier
implementation used $[\SI{0.5}{\newton},\,2\mB g=\SI{40.5}{\newton}]$, a box
derived from the vehicle's \emph{weight} rather than its actuators, and therefore
over-stated the available thrust by \SI{29}{\percent}. Planning against authority
the vehicle does not have is not a conservative error; the surplus is removed
silently by the clip in \eqref{eq:throttleinv}, so the planned recovery
acceleration is never delivered, the tracking error grows, and the next solve
asks for still more thrust.

\section{Estimation: Mass and First Moment as the Online Degrees of Freedom}\label{sec:model:mhe}

\subsection{Augmented state and estimator formulation}\label{sec:model:mheform}

The moving-horizon estimator augments both payload states into the state vector,
\begin{equation}
x_{\mathrm{aug}} = \bigl[\,x;\; \mT;\; s_x;\; s_y\,\bigr] \in \mathbb{R}^{16},
\qquad \dot{\mT} = 0, \quad \dot{\svec}=\bm{0},
\label{eq:xaug}
\end{equation}
holding both constant across the horizon so that a window of measurements
constrains a single value of each rather than a free trajectory. Process noise
$\bm{w} \in \mathbb{R}^{13}$ is admitted on the physical states only; adding
noise to the payload channels would allow them to drift within the window and
would dissolve precisely the constraint that makes the estimate meaningful.
Chapter~\ref{ch:framework} gives the numerical form of the optimisation, its
weights and its bounds; the present section states only what the states mean and
which residual identifies which.

The geometry does \emph{not} appear in the decision vector. The estimator's
runtime parameter vector carries the reconstructed wrench, the previous window's
mass estimate required by \eqref{eq:Afrozen}, and the weak vertical arm:
\begin{equation}
p_{\mathrm{MHE}} = \bigl[\, \Tphys;\; \tauphys;\; \hat m^{-};\; \cdot\;;\; r_z \,\bigr] \in \mathbb{R}^{7}.
\label{eq:pmhe}
\end{equation}
Inside the model, $\cvec$ and $\Jt$ are evaluated from the decision variables by
\eqref{eq:c_from_s} and \eqref{eq:Jdeployed}, so both the translational and the
rotational Jacobians of the payload states are consistent within a single solve.
This is a structural point, not a numerical one: an implementation that computed
$\cvec$ and $\dJ$ \emph{outside} the optimiser from the previous mass estimate
would close an outer fixed-point iteration,
\begin{equation*}
\hat m \;\longrightarrow\; \dJ,\ \cvec
\;\longrightarrow\; \text{attitude model}
\;\longrightarrow\; \hat m ,
\end{equation*}
with positive feedback and no cost term to arbitrate it. That loop was observed
to lock onto self-consistent wrong solutions on light payloads.

\subsection{Which residual identifies what}\label{sec:model:observability}

The two payload states enter through disjoint channels, which is what makes
augmenting the second one nearly free in terms of sensing.

The composite mass enters the translational equation \eqref{eq:trans} as $1/\mT$
and is therefore identified predominantly by the \emph{velocity} residual. It is
not, however, absent from the rotational channel: it appears there through
$\cvec=\svec/\mT$ in \eqref{eq:tcom_s} and through the off-diagonal terms
$B_i=(\mB/\mT)s_i r_z$ of \eqref{eq:ABC}, neither of which is frozen. What
\eqref{eq:Afrozen} removes is therefore not every mass dependence of
$\dot{\bm{\omega}}$, but the one that was demonstrably abused: the $A$ term sits
in the \emph{denominator} $J_{xx}+A$ and carries the amplification
$r_z^2/J_{xx}\approx15$, so inflating $\hat m$ measurably suppresses an angular
residual. The paths that remain are first order and unamplified --- a relative
error in $\mT$ produces an equal relative error in $\tcom$ --- and offer the
optimiser no comparable lever.

The first moment $\svec$ does not appear in \eqref{eq:trans} at all. It enters
only through the trim moment \eqref{eq:tcom_s} and the off-diagonal inertia
\eqref{eq:Jxz}, and is therefore identified exclusively by the
\emph{angular-rate} residual.

Both residuals are components of the same odometry measurement. The second state
consequently costs no additional sensing --- only the two extra decision
variables and the solver time they imply, quantified in
Chapter~\ref{ch:evaluation}. Neither quantity is measured directly: both are
recovered as the window-constant values that reconcile the reconstructed input
$u^{\mathrm{phys}}$ with the observed $\bm v$ and $\bm\omega$ trajectories.

It is worth recording why this separation was not available in the predecessor
design. When the rotational channel is parameterised by $(\mP,\bm r_{xy})$ with
$\bm r_{xy}$ a fixed prior, the quantity the angular-rate residual actually
identifies is the product $\mP\bm r_{xy}$ --- so an error in the assumed
$\bm r_{xy}$ is converted directly into a mass error. Numerically, on the
eccentric-hover condition, halving the assumed lever arm biased the mass estimate
by \SI{+14.4}{\percent} and setting it to zero by \SI{+41.6}{\percent}, whereas a
\SI{\pm33}{\percent} error in $r_z$ changed the mass estimate by
\SI{0.00}{\percent}. Augmenting the product removes the first sensitivity
entirely; the second never existed, which is why $r_z$ may safely remain a weak
prior. The insensitivity is structural rather than fortunate: $r_z$ is
annihilated by the cross product in \eqref{eq:tcom}, so the torque channel
carries no information about it in the first place and no amount of excitation
would identify it. A Cram\'er--Rao bound computed for the eccentric-hover window tells the
same story from the other side: with the rotational channel carrying no mass
information the bound is $\sigma_m=\SI{0.0109}{\kg}$, and with the channels
coupled it is $\sigma_m=\SI{0.0011}{\kg}$, with \SI{94}{\percent} of the
information arriving through rotation.

\subsection{Prior-free initialisation}\label{sec:model:seed}

The estimator requires an initial iterate for the mass channel. Supplying the
calibrated airframe mass would be a prior about the answer; the implementation
instead constructs the seed from measurements.

The simplest form is the hover approximation obtained by setting $a_z=0$ and
$\cos\theta=1$ in \eqref{eq:vbal},
\begin{equation}
m^{(0)}=\operatorname{clip}\!\left(\frac{\Tphys}{g},\; \underline m,\; \overline m\right),
\label{eq:massseed}
\end{equation}
which is prior-free because $\Tphys$ comes from \eqref{eq:phys} and involves no
model state. Its weakness is that it is only valid near hover: on a
\SI{4}{\meter\per\second} figure-eight the instantaneous ratio $\Tphys/g$ swings
from \SI{-42}{\percent} to \SI{+30}{\percent} of the true mass, because the
collective thrust is also supplying centripetal acceleration and is tilted away
from the vertical.

The deployed seed therefore uses the whole window rather than one sample. Summing
the exact vertical balance \eqref{eq:vbal} over the $N$ stages of the window,
\begin{equation}
\mT \sum_{i=0}^{N-1}\bigl(g + a_{z,i}\bigr) \;=\; \sum_{i=0}^{N-1} T_i \cos\theta_i ,
\label{eq:winbal1}
\end{equation}
and eliminating the accelerations by the telescoping identity
\begin{equation}
\sum_{i=0}^{N-1} a_{z,i}
= \sum_{i=0}^{N-1} \frac{v_{z,i+1}-v_{z,i}}{\Delta t}
= \frac{v_{z,N}-v_{z,0}}{\Delta t},
\label{eq:telescope}
\end{equation}
gives a seed that requires no numerical differentiation of the velocity signal at
all:
\begin{equation}
m^{(0)} \;=\; \frac{\displaystyle\sum_{i=0}^{N-1} T_i \cos\theta_i}
                   {\displaystyle N g + \frac{v_{z,N}-v_{z,0}}{\Delta t}} .
\label{eq:winbal}
\end{equation}
Only the two velocity endpoints of the window enter, so the estimate is immune to
the sample-to-sample noise that a finite-difference acceleration would amplify;
the tilt factor $\cos\theta_i$ from \eqref{eq:costheta} restores the vertical
component that \eqref{eq:massseed} discards; and every input is either $\Tphys$
or odometry, so the seed remains decoupled from the estimator's own state. The
implementation falls back to \eqref{eq:massseed}, and then to a neutral box
value, when the window is not yet full or the denominator of \eqref{eq:winbal}
degenerates.

\paragraph{Initial guess versus prior.} The distinction matters for the claims of
this dissertation and is easy to blur. Equation~\eqref{eq:winbal} enters the
solver in two places: the initial iterate of the SQP, and the mean of the arrival
cost. The former affects only the number of iterations, never the optimum. The
latter affects the optimum in proportion to its weight, which is deliberately set
so low that the corresponding standard deviation is several kilograms --- far
wider than the range of masses the vehicle can carry. The seed therefore
influences where the optimiser starts, not where it converges. The same
expression is used on re-initialisation after repeated solver failure, for the
same reason: restarting from a measurement is safer than restarting from a
possibly corrupted previous estimate.

\section{Control: How the Parameters Re-centre the Cost}\label{sec:model:ctrl}

The nonlinear model predictive controller carries the parameters as runtime
values,
\begin{equation}
p_{\mathrm{NMPC}} = \bigl[\, x_{\mathrm{ref}}(13);\; \mT;\; \dJ;\; \cvec(2);\; \bm{d}(3);\; \bm{\xi}(3) \,\bigr] \in \mathbb{R}^{23},
\label{eq:pnmpc}
\end{equation}
where $\bm{d}$ and $\bm{\xi}$ are lumped translational and rotational
disturbances used only by the $\mathcal{L}_1$ comparison baseline
\cite{hanover2021} and identically zero under the proposed method. Over a
prediction horizon of $N$ stages it solves
\begin{align}
\min_{u_{0:N-1}} \quad
& \sum_{k=0}^{N-1} \Bigl( \bigl\| e_{\mathrm{track}}(x_k, x_{\mathrm{ref},k}) \bigr\|^2_{Q}
+ \bigl\| u_k - u_{\mathrm{hover}}(p) \bigr\|^2_{R} \Bigr) \nonumber \\
& \qquad + \bigl\| e_{\mathrm{track}}(x_N, x_{\mathrm{ref},N}) \bigr\|^2_{P} \nonumber \\
\mathrm{s.t.} \quad
& x_{k+1} = F(x_k, u_k;\, p_{\mathrm{NMPC}}), \nonumber \\
& \underline{T} \le T_k \le \overline{T}, \quad
  |\tau_{k,i}| \le \overline{\tau}_i ,
\label{eq:nmpc}
\end{align}
applying only $u_0$ and re-solving at the next step, with the thrust box taken
from the actuator envelope \eqref{eq:envelope}.

\subsection{Why the input baseline must track the parameters}\label{sec:model:hover}

The second term of \eqref{eq:nmpc} penalises the deviation of the input from a
\emph{baseline} $u_{\mathrm{hover}}$, and that baseline is recomputed from the
live parameters at every solve rather than fixed when the solver is generated:
\begin{equation}
T_h = \mT\bigl(g - d_z\bigr),
\qquad
u_{\mathrm{hover}} = \bigl[\, T_h,\; c_y T_h,\; -c_x T_h,\; 0 \,\bigr]^{\!\top}.
\label{eq:hover}
\end{equation}

Consider first why a baseline is needed at all. Penalising $\|u\|^2_R$ outright
would treat the thrust required merely to remain airborne as an expense to be
minimised, biasing the solution toward insufficient thrust. The baseline
recentres the penalty on the input that sustains equilibrium, so that $R$
penalises only genuine control effort.

Consider next what happens when the baseline is \emph{stale}. Suppose
$u_{\mathrm{hover}}$ is computed from the unloaded mass while the vehicle carries
a parcel. The cost then pulls the thrust persistently toward a value below what
equilibrium requires, and the optimum settles where the marginal tracking cost
balances the marginal baseline-deviation cost --- at a nonzero steady-state
error. The behaviour has a counter-intuitive signature that identifies it in
practice: \emph{increasing $R$ makes the error worse}, because $R$ is precisely
the strength with which the incorrect baseline pulls. No amount of weight tuning
removes a bias of this kind; only correcting the baseline does. The effect is not
marginal --- with a heavy load and an empty-airframe baseline the tracking error
was observed to stall above \SI{0.85}{\meter} without converging.

The torque entries of \eqref{eq:hover} are exactly $-\tcom$ from \eqref{eq:tcom},
and they exist for the same reason. Hovering with an eccentric load requires a
constant trim torque (Section~\ref{sec:model:tcom}); a cost that does not admit
this pulls the torque toward zero and produces the same class of steady-state
bias in the attitude channel. The two equations are complementary and neither
suffices alone: \eqref{eq:tcom} teaches the \emph{model} that the moment exists,
while \eqref{eq:hover} teaches the \emph{objective} that it is necessary. The
$d_z$ term in $T_h$ plays the identical role for the lumped-disturbance baseline:
with $\bm d$ in the model, vertical equilibrium is $T/\mT-g+d_z=0$, so omitting
$d_z$ from the penalty centre would reintroduce the very bias the baseline exists
to remove.

\subsection{Inner-loop bandwidth}\label{sec:model:gain}

Model consistency in the outer loop cannot restore inner-loop bandwidth. The rate
controller running on the flight controller is tuned for the bare airframe; a
payload that multiplies the roll and pitch inertia lowers its effective
bandwidth, and for heavy eccentric loads the resulting attitude oscillation grows
until the torque saturates. The relevant ratio is
\begin{equation}
\kappa \;=\; \frac{J_{xx} + \dJ}{J_{xx}},
\label{eq:gain}
\end{equation}
which \eqref{eq:numbers} puts at $5.30$ for the heaviest payload considered: that
payload sits at the ceiling of what inner-loop retuning alone can absorb on this
airframe. Together with the torque budget \eqref{eq:trimtorque}, this defines a
feasibility boundary on payload mass and eccentricity that is characterised
experimentally in Chapter~\ref{ch:evaluation}.

The predecessor implementation applied $\kappa$ by rewriting the flight
controller's rate-loop gains at attachment. The deployed implementation does not
touch those parameters; it scales the body-rate \emph{command} instead, which
keeps the adaptation inside the ROS layer, makes it continuous, and makes it
reversible. The mechanism is derived in Section~\ref{sec:frame:omegascale}.

\section{What is Estimated, What is Derived, and What is Not}\label{sec:model:whynot}

Table~\ref{tab:params} records the identity of every parameter in the model. Two
scalars are online degrees of freedom; everything else is either derived from
them by the algebra of this chapter, calibrated once for the airframe, or
retained as a weak prior whose error has been shown not to matter.

\begin{table}[ht]
\caption{Model parameters, classified by identity. Only the composite mass and
the two components of the horizontal first mass moment are decision variables of
the estimator; the inertia increment and the CoM offset are generated from them.}
\label{tab:params}
\centering
\footnotesize
\begin{tabular}{@{}l p{0.20\textwidth} c p{0.26\textwidth} p{0.16\textwidth}@{}}
\toprule
Symbol & Meaning & Dim. & Identity & Enters through \\
\midrule
$\mT$ & composite mass & 1 & \textbf{MHE decision variable} & \eqref{eq:trans}, \eqref{eq:hover} \\
$\svec$ & horizontal first mass moment & 2 & \textbf{MHE decision variable} & \eqref{eq:tcom_s}, \eqref{eq:Jxz} \\
\midrule
$\cvec$ & CoM horizontal offset & 2 & derived, $\svec/\mT$ & \eqref{eq:tcom}, \eqref{eq:hover} \\
$\dJ$ & roll/pitch inertia increment & 1 & derived, $A$ of \eqref{eq:Jdeployed} & \eqref{eq:rot}, \eqref{eq:gain} \\
$\mP$ & payload mass & 1 & derived, $(\mT-\mB)^{+}$ & \eqref{eq:Jdeployed} \\
\midrule
$r_z$ & payload vertical arm & 1 & weak prior, \SI{-0.47}{\meter} & \eqref{eq:Jdeployed}, \eqref{eq:Jxz} \\
$\rxy$ & payload horizontal arm & 2 & \textbf{not used} & --- \\
\midrule
$\mB$ & bare-airframe mass & 1 & calibration & \SI{2.0643}{\kg} \\
$\Jb$ & bare-airframe inertia & 3 & calibration & $\mathrm{diag}(0.0142,$ $0.0142, 0.0210)$ \\
$k_f, k_m$ & thrust and yaw-reaction coefficients & 2 & calibration & \SI{8.55e-6}{}, $0.016$ \\
$k_d$, $g$ & drag coefficient, gravity & 2 & constants & $0.05$, $9.81$ \\
\bottomrule
\end{tabular}
\end{table}

Three design choices in that table invite objection, and each is answered by an
argument already made above.

\paragraph{Why not augment the inertia increment as well?} Because it is nearly
unobservable in the flight regime of interest. The increment acts only through
$\dot{\bm{\omega}}$ in \eqref{eq:rot}, and a delivery task --- spent largely in
near-hover, low-rate flight --- excites that channel weakly. Admitting $\dJ$ as a
free variable would let the optimiser explain measurement noise with a quantity
the data barely constrains. There is a second, sharper reason: $\Jt$ appears
\emph{inverted} in \eqref{eq:rot} and is propagated $N$ steps forward, so an
error in it is amplified rather than damped, whereas the additive lumped term
$\bm\xi$ of \eqref{eq:pnmpc} --- the alternative used by disturbance-based
methods --- is always safe. Estimating six inertia elements would also add
decision variables whose natural scale is $10^{-2}$, two orders of magnitude
below the mass, to a Hessian already sensitive to scaling.

\paragraph{Why not augment the horizontal arm?} Because nothing observes it.
Every appearance of $\rxy$ in the equations of motion is inside the
product $\mP\rxy$, which is the state $\svec$ we do augment. Estimating
the factors separately would add an unidentifiable direction, and would
reintroduce the release singularity of Section~\ref{sec:model:whyS}.

\paragraph{Why is a prior on $r_z$ acceptable when a prior on $\rxy$ is
not?} Because their error sensitivities differ by orders of magnitude. An error
in $\rxy$ propagates directly into the mass estimate through the product
the rotational residual identifies; an error in $r_z$ enters only $A$ and the
off-diagonal terms, both of which are second-order effects in near-hover flight
where $\dot{\bm{\omega}}\approx\bm 0$. A \SI{\pm33}{\percent} perturbation of
$r_z$ moved the mass estimate by less than the reporting precision.

\paragraph{The information boundary.} One mass-valued constant does remain in the
closed loop, and it is disclosed rather than absorbed: the platform's
\SI{0.30}{\kg} payload \emph{envelope}. It is used for the attach-time inertia
bootstrap of Section~\ref{sec:frame:bootstrap}, for the physical admissibility
bound on $\svec$ in Section~\ref{sec:frame:momentref}, and for the envelope
treatment of Table~\ref{tab:amode}. It is a specification of the airframe, in the
same category as $\Jb$ or the actuator envelope \eqref{eq:envelope}, and it is
not the mass of the parcel being carried on any particular flight. The precise
claim of this dissertation is therefore: \emph{no task-specific payload mass, and
no attach or release notification, enters the estimator or the controller's model
update}; it is not the stronger claim that the flight stack contains no
payload-scale constant whatsoever.

\section{Declared Modelling Approximations}\label{sec:model:approx}

Four approximations are stated explicitly here rather than left to be discovered.

\paragraph{The second-order horizontal inertia terms are discarded.} The group
$C_{ij}$ of \eqref{eq:ABC} is set to zero, as derived in
Section~\ref{sec:model:Jofs}. The error is \eqref{eq:Cerror},
\SI{2.6}{\percent} of the roll/pitch increment at the platform geometry, and the
reason is the degeneracy \eqref{eq:Cdegenerate}. This is a deliberate exchange of
a small bias for a well-posed release limit.

\paragraph{Anisotropy of the inertia increment is not represented.} The deployed
model \eqref{eq:Jdeployed} carries a single $\dJ$ on both roll and pitch, whereas
\eqref{eq:djaxes} gives $\dJ_{xx}\ne\dJ_{yy}$ whenever $r_x\ne r_y$. After the
$C_{ij}$ terms are discarded the two become identical, $\dJ_{xx}=\dJ_{yy}=A$, so
this approximation is subsumed by the previous one rather than additional to it.

\paragraph{The yaw increment is neglected.} The third expression in
\eqref{eq:djaxes}, $\dJ_{zz} = \mu(r_x^2 + r_y^2)$, is not added to $J_{zz}$. At
\SI{0.3}{\kg} and $r_{xy} = \SI{0.109}{\meter}$ it amounts to
\SI{0.0031}{\kg\meter\squared}, or \SI{14.8}{\percent} of $J_{zz}$. This is a
known and non-negligible approximation. It is retained because in the deployed
coordinates $\dJ_{zz}$ belongs to the $C_{ij}$ group and would reintroduce
\eqref{eq:Cdegenerate}, and because the yaw channel neither carries payload
compensation nor participates in any release decision.

\paragraph{The vertical CoM offset is absent by construction.} As shown in
Section~\ref{sec:model:tcom}, $c_z$ drops out of \eqref{eq:tcom} as an exact
algebraic consequence of collinearity, not as a truncation. It is listed among
the approximations only to forestall the opposite reading.

\paragraph{The inertia is treated as constant in the body frame.}
Equation~\eqref{eq:rot} omits the $\dot{\Jt}\bm{\omega}$ term that a
time-varying inertia would contribute. The omission is exact on either side of a
grasp or release and wrong only across one, for the duration of the transfer
itself. Within the estimator the affected samples fall inside a single window and
are absorbed as process noise; no attempt is made to model the transfer
dynamically.

Two further simplifications are inherited from the rotor model rather than the
rigid-body model. The thrust coefficient $k_f$ is a static constant, and is known
to lose validity under large induced inflow during fast translation
\cite{bangura2012}; and the yaw reaction of \eqref{eq:phys} is the steady-state
term only, omitting the rotor angular-acceleration contribution
$\sum_i J_r \dot\omega_i$, which is negligible in hover and mild manoeuvre but
not during rapid yaw. Both are re-examined in Chapter~\ref{ch:evaluation}, where
the sensitivity of the complete interface to a mis-calibrated $k_f$ is measured
directly.

\end{spacing}
\newpage
